%% AUTO-GENERATED by scripts/make_tables.py -- DO NOT EDIT BY HAND.
%% Source: results/comparison.json
%% Status: STUB - no measurements available
\begin{table}[htbp]
\centering\small
\caption{층위 2: 기하를 표면에 고정하는 것의 대가와 이득. 광도 지표와 기하 지표를
\emph{반드시} 함께 읽어야 한다. 자유 기하는 PSNR에서 이기고 $\Esurf$에서 크게
지도록 설계상 예상되며, 그 상충이 본 논문의 논점이다.}
\label{tab:layer2-free}
\begin{tabular}{@{}lrrrrr@{}}
\toprule
표현 & PSNR$_{\mathrm{ROI}}$ (dB) & $\Esurf$ (mm) & 법선 오차 ($^\circ$)
  & 프레임당 바이트 & 기하 고정 \\
\midrule
CV-Dyn2DGS (2D 원반) & \NM
  & \NM & \NM
  & \NM & 예 \\
Gaussian surfel (z-scale $\to0$) & \NM
  & \NM
  & \NM
  & \NM & 예 \\
thin-3DGS ($\rho=0.1$) & \NM
  & \NM & \NM
  & \NM & 예 \\
자유 3DGS (표면 초기화) & \NM
  & \NM
  & \NM
  & \NM & \textbf{아니오} \\
자유 3DGS (무작위 초기화) & \NM
  & \NM
  & \NM
  & \NM & \textbf{아니오} \\
\bottomrule
\end{tabular}
\end{table}
